\subsection{Potenziale a gradino}
\begin{figure}
\centering
\begin{tikzpicture}
  \draw (-5,0) -- (5,0);
  \draw (0,-1) -- (0,3);

  \draw[thick] (-4, 0) -- (0,0);
  \draw[thick] (0,2) -- (4,2);

  \node (1) at (-2.5,2.5) {1};
  \node (2) at (2.5,2.5) {2};
  \node (V0) at (-.3,2) {$V_0$};
\end{tikzpicture}
\caption{Schema del potenziale a gradino.}
\label{gradino}
\end{figure}
Lo schema \`e riportato in Figura \ref{gradino}.
$$
V(x) = \left\{ \begin{matrix}
V_0 & x > 0 \\
0 & x < 0
\end{matrix} \right..
$$
Supponendo che
$$
\delta(x) = \frac{1}{V_0} \frac{d}{dx} V(x),
$$
perch\'e
\begin{align*}
\int_{-\infty}^{+\infty} V(x) f'(x) \, dx &= \left[ V(x) f(x) \right]_{-\infty}^{+\infty} - \int_{-\infty}^{+\infty} f(x) \frac{d}{dx}V(x) \, dx = \\
&= f(+\infty) - \int_{-\infty}^{+\infty} f(x) \frac{d}{dx}V(x) \, dx \\
\int_{-\infty}^{+\infty} V(x) f'(x) \, dx &= \int_0^{+\infty} f'(x) \, dx = f(+\infty) - f(0).
\end{align*}

Scrivo l'equazione di Schr\"odinger
$$
-\frac{\hbar^2}{2m} \psi''(x) + V(x) \psi(x) = E \psi(x)
$$
$$
\psi''(x) = -\frac{2m}{\hbar^2} (E - V(x)) \psi(x).
$$
Allora la $ \psi'' $ \`e discontinua perch\'e lo \`e $ V $, ma quindi $ \psi $ e $ \psi' $ sono continue, cos\`{i} posso dividere il problema in due parti e poi imporre delle condizioni di raccordo.

Per $ x < 0 $ ho $ \psi_1 $ tale che 
$$
\psi_1''(x) = -\frac{2m}{\hbar^2} E \psi_1(x) \; \Rightarrow \; \psi_1(x) = A e^{ikx} + B e^{-ikx}, \; k = \frac{\sqrt{2mE}}{\hbar}.
$$

Per $ x > 0 $
\begin{align*}
\psi_2''(x) &= -\frac{2m}{\hbar^2} (E - V_0) \psi_2(x) \; \Rightarrow \\
&\; \Rightarrow \; \psi_2(x) = C e^{ik'x} + D e^{-ik'x}, \; k' = \sqrt{\frac{2m}{\hbar^2}(E - V_0)}
\end{align*}
e per ora impongo $ E > V_0 $.

Poich\'e $ \psi_1(0) = \psi_2(0) $ e $ \psi_1'(0) = \psi_2'(0) $, allora mi restano due condizioni da imporre. Sfrutto la normalizzazione per la terza equazione e poi impongo $ D = 0 $. Allora 
$$
\begin{cases}
A + B = C \\
ik(A - B) = ik'C 
\end{cases} \quad \Rightarrow \quad \begin{cases}
B = \frac{k - k'}{k + k'} A \\
C = \frac{2k}{k + k'} A
\end{cases}.
$$

Introduco la \textbf{corrente di probabilit\`a}
$$
J(x, t) = -\frac{i\hbar}{2m} \left( \psi^*(x, t) \frac{\partial}{\partial x} \psi(x, t) - \left( \frac{\partial}{\partial x} \psi^*(x, t) \right) \psi(x, t) \right)
$$
tale che
$$
\frac{\partial}{\partial x} J(x, t) + \frac{\partial}{\partial t} \rho(x, t) = 0, \quad \text{con } \rho(x, t) = |\psi(x, t)|^2 = |\braket{x|\psi(t)}|^2.
$$
Infatti
$$
\frac{\partial J}{\partial x} = -\frac{i\hbar}{2m} \left( \cancel{\frac{\partial \psi^*}{\partial x} \frac{\partial \psi}{\partial x}} + \psi^* \frac{\partial^2 \psi}{\partial x^2} - \frac{\partial^2 \psi^*}{\partial x^2} \psi - \cancel{\frac{\partial \psi^*}{\partial x} \frac{\partial \psi}{\partial x}} \right)
$$
poi per l'equazione di Schr\"odinger
$$
\frac{\partial^2 \psi}{\partial x^2} = \left( i\hbar\frac{\partial \psi}{\partial t} -V(x)\psi \right) \left( -\frac{2m}{\hbar^2} \right).
$$

Essendo $ \mathcal{H} $ hermitiano, $ \mathcal{H} = \mathcal{H}^\dag $, quindi 
$$
\frac{\partial^2 \psi^*}{\partial x^2} = \left( -i\hbar\frac{\partial \psi^*}{\partial t} - V(x) \psi^* \right) \left( -\frac{2m}{\hbar^2} \right)
$$
allora 
\begin{align*}
\frac{\partial}{\partial x} J(x, t) &= -\frac{i\hbar}{2m} \left( \psi^* \left( i\hbar\frac{\partial \psi}{\partial t} - V \psi \right) - \left( -i\hbar\frac{\partial \psi^*}{\partial t} - V \psi^* \right) \psi \right) \left( -\frac{2m}{\hbar^2} \right) = \\
&= - \left( \psi^* \frac{\partial \psi}{\partial t} + \frac{\partial \psi^*}{\partial t} \psi \right) = - \frac{\partial}{\partial t} \left( \psi^* \psi \right) = \\
&= -\frac{\partial}{\partial t} \rho(x, t).
\end{align*}

Integrando
$$
\int_a^b \frac{\partial \rho}{\partial t} \, dx = -\int_a^b \frac{\partial J}{ \partial x} \, dx = - \left( J(b) - J(a) \right)
$$
$$
\frac{\partial \rho}{\partial t} > 0 \qquad \Rightarrow \qquad J(a) > J(b).
$$

Il valore medio di $ J $
\begin{align*}
\int_{-\infty}^{+\infty} J(x, t) \, dx &= \frac{1}{m} \int_{-\infty}^{+\infty} \braket{\psi|x} \braket{x|\hat{p}|\psi} \, dx = \\
&= \frac{1}{m} \braket{\psi|\hat{p}|\psi} = \\
&= \frac{\braket{\hat{p}}_\psi}{m}
\end{align*}
corrispondente al valore medio di $ \hat{p} $ diviso per $ m $
$$
\braket{J} = \frac{1}{m} \braket{k|\hat{p}|k} = \frac{k}{m} \braket{k|k}.
$$
Quindi data
$$
\psi(x, t) = A \exp \left\{ \frac{i}{\hbar} (kx - \omega t) \right\} + B \exp \left\{ -\frac{i}{\hbar} (kx - \omega t) \right\}
$$
la $ J $ associata \`e 
\begin{align*}
J(x, t) &= -\frac{i\hbar}{2m} \left( \left( A^* \exp \left\{ -\frac{i}{\hbar} (kx - \omega t) \right\} + B^* \exp \left\{ \frac{i}{\hbar} (kx - \omega t) \right\} \right) \right. \cdot \\
&\cdot \frac{ik}{\hbar} \left( A \exp \left\{ \frac{i}{\hbar} (\dots) \right\} + B \exp \left\{ -\frac{i}{\hbar} (\dots) \right\} \right) + \\
&- \frac{ik}{\hbar} \left( -A^* \exp \left\{ -\frac{i}{\hbar} (\dots) \right\} + B^* \exp \left\{ -\frac{i}{\hbar} (\dots) \right\}  \right) \cdot \\
&\cdot \left( A \exp \left\{ \frac{i}{\hbar} (\dots) \right\} + B \exp \left\{ -\frac{i}{\hbar} (\dots) \right\} \right) = \\
&= \frac{k}{2m} \left( |A|^2 - |B|^2 \right) 2 = \frac{k}{m} \left( |A|^2 - |B|^2 \right) = \\
&= J_A(x, t) + J_B(x, t).
\end{align*}

Poich\'e $ \ket{\psi} = \ket{k} \braket{k|\psi} + \ket{-k} \braket{-k|\psi} $, allora 
$$
\braket{J} = \frac{k}{m} \left( |\braket{k|\psi}|^2 - |\braket{-k|\psi}|^2 \right).
$$
Inoltre, essendo 
$$
\rho = |\psi|^2 \; \Rightarrow \; \frac{\partial \rho}{\partial t} = 0 \; \Rightarrow \; \frac{\partial J}{\partial x} = 0,
$$
quindi $ J $ \`e invariante per spostamento, quindi 
$$
J_A + J_B = J_C \text{ costante}
$$
e questa condizione deve valere nella prima e nella seconda regione su cui poi impongo le equazioni di continuit\`a. Allora 
$$
J_{\text{inc.}} + J_{\text{rifl.}} = J_{\text{trasm.}}
$$
e i coefficienti
$$
T = \left| \frac{C}{A} \right|^2, \qquad R = \left| \frac{B}{A} \right|^2.
$$

Nel caso in cui $ E < V_0 $
$$
k' = \sqrt{\frac{2m}{\hbar}(E - V_0)} \quad \Rightarrow \quad k' = i \beta, \, \beta \in \mathbb{R}
$$
quindi $ \psi_2(x) = C e^{-\beta x} $, cio\`e l'andamento \`e esponenziale e la corrente $ J_C(x, t) = 0 $, quindi $ J_A = - J_B $ e 
$$
R = \left| \frac{B}{A} \right|^2 = \left| \frac{k - i\beta}{k + i\beta} \right|^2 = 1.
$$
Se $ V_0 \to +\infty $, allora $ \beta \to +\infty $, allora $ \psi_2(x) = 0 $, quindi essendo $ \psi_1(0) = \psi_2(0) $, $ A + B = 0 $, $ \psi_1(x) = A' (e^{ikx} - e^{-ikx}) $.